% ===========================================================================
\subsection{\file{new-design-reference/js/main.js}}
% ===========================================================================

\textbf{JavaScript, 64 lines.} Everything on the page that reacts. It does three
things and nothing else: the navigation bar grows a frosted background once you
scroll, the four metric numbers count up from zero when they come into view, and
most blocks fade upward as they arrive.

\begin{provenance}{new-design-reference/js/main.js}
Hand-written, and its own header comment says as much on line 2: \code{Zero
external dependencies.} That claim is true and checkable: the file imports nothing,
requires nothing, and the repository has no \file{package.json} that could supply a
library.

It arrived complete in the first commit \code{1f7ed12} and has never been edited.
It is a condensed version of the sibling repository's file: same three behaviours,
same helper, same counter arithmetic, 15 lines shorter, with the comments trimmed
and one behaviour dropped (the sibling watches a fifth selector this page does not
have). Which was written first is not recorded and is not guessed here.
\textbf{Confidence: certain for hand-written and dependency-free; the relationship
to the sibling is stated as a close resemblance, not as a direction of copying.}
\end{provenance}

\subsubsection*{Why it exists}

Without this file the page is complete, readable and static. Two of its three jobs
are pure decoration and would not be missed. The third matters: the metric numbers
are written into the HTML as the character \code{0} with their real values hidden
in data attributes, so with this file gone the metrics band permanently advertises
zero of everything.

That damage is smaller here than in the sibling, because the three hero statistics
(\file{index.html} lines 114 to 127) state three of the same four numbers as plain
text. A visitor without JavaScript still sees ``17+ projects'', ``1.6M+ organic
views'' and ``10K+ users served'' in the hero, and four zeros further down. Odd,
but not silent.

\begin{keyidea}{The single most important line in this file is in another file}
\file{index.html} loads this script at line 393, immediately before
\code{</body>}. Everything below runs the instant it loads, with no wrapper telling
it to wait for the document. That only works because by then every element already
exists. Move the \code{<script>} tag into the \code{<head>} and this file would run
against an empty page, find nothing, attach nothing, and report no error at all.
\end{keyidea}

\subsubsection*{Line by line}

\paragraph{Lines 1 to 13: the header, the wrapper, and the navigation bar.}

\begin{lstlisting}[language=Java,firstnumber=1,caption={\file{js/main.js}, lines 1--13}]
/* salmanadnan.com · main.js
   Scroll-driven animations. Zero external dependencies. */

(function () {
  'use strict';

  /* ── Nav gloss on scroll ── */
  const nav = document.querySelector('.nav');
  if (nav) {
    const tick = () => nav.classList.toggle('scrolled', window.scrollY > 40);
    window.addEventListener('scroll', tick, { passive: true });
    tick();
  }
\end{lstlisting}

Lines 1 and 2 are a comment, using the same \code{/* */} form as CSS.

Line 4 opens a wrapper that looks like noise and is not.

\begin{jargon}{Immediately invoked function expression}
Writing \code{(function () \{ ... \})();} defines a function and runs it once,
straight away, without giving it a name. The point is not the running: it is the
boundary. Anything declared inside is invisible outside.

Without it, the names \code{nav}, \code{animateCounter} and \code{observe} declared
below would become global, visible to any other script on the page, and any other
script using the same name would collide with them. One stray global called
\code{observe} is exactly the kind of bug that takes an afternoon to find.
\end{jargon}

Line 5 is \code{'use strict';}, switching the language into a stricter mode for
everything inside the wrapper.

\begin{jargon}{Strict mode}
A setting making JavaScript refuse things it would otherwise silently allow. The
most valuable one: without it, assigning to a name you never declared quietly
creates a global variable, so a typo such as \code{navv = 3} succeeds and does
nothing useful. In strict mode it throws an error and you find the typo
immediately.
\end{jargon}

Line 8 finds the navigation bar. \code{document.querySelector} takes a CSS selector,
exactly the syntax the stylesheet uses, and returns the first matching element or
the special value \code{null} if there is none.

\begin{jargon}{const, and why not var}
\code{const} declares a name whose binding cannot later be pointed at something
else. It is the default choice in modern JavaScript, with \code{let} for names that
genuinely must change. The older \code{var} has confusing scoping rules and appears
nowhere in this file, which is a small sign of care.
\end{jargon}

Line 9 is a guard: everything happens only \code{if (nav)}, that is, only if
something was found. This is the pattern making the whole file tolerant of a
changing page.

\begin{jargon}{Truthy and falsy}
JavaScript treats certain values as false in a condition: \code{false}, \code{0}, an
empty string, \code{null} and \code{undefined}. Everything else counts as true. So
\code{if (nav)} reads as ``if an element was found'', because
\code{querySelector} returns \code{null} when it finds nothing.
\end{jargon}

Line 10 does in one line what the sibling spreads over three, using an arrow
function with an implicit return.

\begin{jargon}{Arrow function}
\code{() => expression} is a compact way to write a function. Written without
braces, as here, it returns the expression's value automatically. The return value
is discarded in this case; the compactness is the only benefit.
\end{jargon}

\code{classList.toggle('scrolled', condition)} adds the class when the condition is
true and removes it when false. The condition is \code{window.scrollY > 40}, meaning
the page has scrolled more than 40 pixels. The stylesheet at lines 101 to 107
defines what \code{scrolled} looks like, and the transition on line 99 makes the
change fade rather than snap.

Line 11 registers the function to run on every scroll event. \code{\{ passive: true
\}} is a performance promise to the browser: this handler will never cancel the
scroll. Without the promise the browser must wait for the handler to finish before
knowing whether the scroll may proceed, which can make scrolling feel sticky. It is
the correct flag here and it is often forgotten.

Line 12 calls the function once immediately, handling the case where the page loads
already scrolled: a refresh halfway down, or an arrival on a link ending in
\code{\#book}. Without it the bar would be transparent over content until the first
scroll.

\paragraph{Lines 14 to 33: the counter.}

\begin{lstlisting}[language=Java,firstnumber=14,caption={\file{js/main.js}, lines 14--33}]

  /* ── Counter animation ── */
  function animateCounter(el) {
    const raw    = el.dataset.target || '0';
    const suffix = el.dataset.suffix || '';
    const target = parseFloat(raw.replace(/[^0-9.]/g, ''));
    const isFloat = raw.includes('.');
    const dec     = isFloat ? (raw.split('.')[1] || '').length : 0;
    const dur     = 2200;
    const start   = performance.now();

    const step = (now) => {
      const p     = Math.min((now - start) / dur, 1);
      const eased = 1 - Math.pow(1 - p, 3);
      const val   = eased * target;
      el.textContent = (isFloat ? val.toFixed(dec) : Math.round(val).toLocaleString()) + suffix;
      if (p < 1) requestAnimationFrame(step);
    };
    requestAnimationFrame(step);
  }
\end{lstlisting}

This is the most substantial piece of logic in the repository and is worth taking
slowly.

Line 17 reads the target. \code{el.dataset.target} is how JavaScript reads the
attribute \code{data-target} written in the HTML: strip the \code{data-} prefix and
what remains is a property of \code{dataset}. The \code{|| '0'} supplies a default:
if the attribute is missing, \code{dataset.target} is \code{undefined}, which is
falsy, so the expression evaluates to \code{'0'}.

Line 18 does the same for the suffix, defaulting to an empty string.

Line 19 turns the text into a number. \code{raw.replace(/[\^{}0-9.]/g, '')} deletes
every character that is not a digit or a dot, and \code{parseFloat} reads what is
left as a number.

\begin{jargon}{Regular expression}
A pattern for matching text, written between slashes. \code{[0-9.]} means ``a digit
or a dot''; a leading caret negates it, so it means ``anything that is not a digit
or a dot''. The \code{g} after the closing slash means global: replace every
occurrence, not just the first.
\end{jargon}

Line 20 asks whether the original text contained a dot, which is how the code
distinguishes \code{1.6} from \code{17}. Line 21 counts the digits after the dot so
\code{1.6} displays to one decimal place throughout rather than as
\code{1.5999999}. \code{(raw.split('.')[1] || '')} takes everything after the dot,
falling back to an empty string, and \code{.length} counts it.

\begin{jargon}{Ternary operator}
\code{condition ? a : b} evaluates to \code{a} when the condition is true and
\code{b} otherwise. It is an \code{if} statement compressed into an expression, and
it appears twice in this file, on lines 21 and 29.
\end{jargon}

Line 22 fixes the duration at 2200 milliseconds. Line 23 records the start time with
\code{performance.now()}, a clock designed for measuring durations precisely.

Lines 25 to 31 define one frame, and this is the heart of it.

\code{const p = Math.min((now - start) / dur, 1)} computes progress from 0 to 1. The
\code{Math.min(..., 1)} caps it, so if the browser is busy and a frame arrives late,
progress never exceeds 1 and the number never overshoots. That cap is the difference
between a counter landing on 17 and one that occasionally flashes 19 before
settling.

\code{const eased = 1 - Math.pow(1 - p, 3)} is the easing. Linear progress feels
mechanical; this curve starts fast and slows dramatically near the target, which is
what makes it feel like settling.

\begin{keyidea}{What the easing formula actually does}
Take progress \code{p}, subtract it from 1, cube the result, subtract that from 1.
At \code{p = 0.5}, halfway through in time, the formula gives $1 - 0.5^3 = 0.875$:
the counter is already 87.5\% of the way there. The whole second half of the
animation covers the last eighth of the distance. That asymmetry is the entire
effect.
\end{keyidea}

Line 29 renders the number. If the target had a decimal point, \code{toFixed(dec)}
formats it to that many places. Otherwise \code{Math.round} makes it whole and
\code{toLocaleString()} inserts the thousands separators appropriate to the
visitor's region, so a large number reads as \code{1,200} for a British visitor and
\code{1.200} for a German one. The suffix is glued on the end.

Line 30 asks for the next frame if the animation is unfinished, and line 32 starts
the first.

\begin{jargon}{requestAnimationFrame}
A request to run a function just before the browser's next screen repaint, about 60
times a second. Calling it repeatedly from inside itself, as here, is the standard
way to run an animation: it stays in step with the display, and the browser pauses
it entirely when the tab is in the background, which a timer would not.
\end{jargon}

\begin{honest}{The counter ignores the reduced-motion preference}
This function animates for 2.2 seconds regardless of whether the visitor has asked
their operating system to reduce motion. The stylesheet honours that preference with
unusual care, going further than the sibling by explicitly stopping the floating
badges, so the intent is plainly there and simply was not carried into this file.

The fix is small, and this document states it precisely so nobody has to design it:

\begin{lstlisting}[language=Java,caption={A four-line fix. This code is a
suggestion in this document and is NOT present in the repository.}]
const reduce = window.matchMedia('(prefers-reduced-motion: reduce)').matches;
if (reduce) {
  el.textContent = (isFloat ? target.toFixed(dec)
                            : target.toLocaleString()) + suffix;
  return;
}
\end{lstlisting}

Placed after line 23, that shows the final number immediately and skips the
animation, which is the correct behaviour: the information is preserved and only the
movement is removed.
\end{honest}

\paragraph{Lines 34 to 48: the observer helper.}

\begin{lstlisting}[language=Java,firstnumber=34,caption={\file{js/main.js}, lines 34--48}]

  /* ── Intersection observer helper ── */
  function observe(selector, className, threshold, cb) {
    const els = document.querySelectorAll(selector);
    if (!els.length) return;
    const io = new IntersectionObserver((entries) => {
      entries.forEach((e) => {
        if (!e.isIntersecting) return;
        e.target.classList.add(className || 'visible');
        if (cb) cb(e.target);
        io.unobserve(e.target);
      });
    }, { threshold: threshold || 0.18 });
    els.forEach((el) => io.observe(el));
  }
\end{lstlisting}

This is the best-designed thing in the repository: a small general-purpose helper
that the three lines following it reuse for three different jobs.

\begin{jargon}{IntersectionObserver}
A browser facility that watches elements and reports when they enter or leave the
visible area. Before it existed, the only way to do this was to listen for every
scroll event and measure each element's position by hand, which ran hundreds of
times a second and made pages stutter. The observer does the work outside the main
thread and reports only when something changes.
\end{jargon}

Line 37 finds every matching element. \code{querySelectorAll} returns a list, which
may be empty, rather than \code{null}. Line 38 guards against an empty list and
returns early: \code{!els.length} reads as ``if the count is zero'', because
\code{0} is falsy and \code{!} inverts it.

Lines 39 to 46 create the observer. Its callback receives one entry per element
whose visibility changed, and for each it skips anything leaving rather than
entering, adds the class (defaulting to \code{visible}), runs the optional extra
callback, and stops watching that element.

That last step is the one that matters most.

\begin{keyidea}{Why it stops watching}
Every reveal in this design is meant to happen once. Without \code{unobserve}, an
element scrolling out of view and back would fire again, and for the metrics that
would mean the numbers restarting their count every time the visitor scrolled past.
Unobserving also releases the browser's bookkeeping, so what starts as a dozen
watched elements steadily becomes none.
\end{keyidea}

Line 46 sets the threshold, defaulting to 0.18: the callback fires once 18\% of the
element is visible rather than at the first pixel. Waiting for a fraction is what
makes the reveal feel connected to the scroll instead of triggering just off screen.

\begin{honest}{The defaults cannot express a falsy value}
\code{threshold || 0.18} substitutes the default whenever the argument is falsy, and
\code{0} is falsy. A caller asking for a threshold of \code{0}, meaning ``fire the
moment a single pixel appears'', would silently get \code{0.18}. The same objection
applies to \code{className || 'visible'}, though an empty class name is not a
sensible request.

No caller passes \code{0}, so nothing is broken today. It is a latent trap worth
recognising, and the modern remedy is the \code{??} operator, which substitutes only
for \code{null} and \code{undefined} and leaves \code{0} alone.
\end{honest}

\paragraph{Lines 49 to 52: the three things being watched.}

\begin{lstlisting}[language=Java,firstnumber=49,caption={\file{js/main.js}, lines 49--52}]

  observe('.reveal', 'visible', 0.15);
  observe('.stagger', 'visible', 0.15);
  observe('.metric-num[data-target]', 'counted', 0.5, animateCounter);
\end{lstlisting}

Three calls, and the page's entire reveal behaviour. The first two add the class
\code{visible} to anything carrying \code{reveal} or \code{stagger}, both at a 0.15
threshold, and the stylesheet does everything after that.

The third is different. Its selector \code{.metric-num[data-target]} matches
elements carrying the class \emph{and} the attribute, a small piece of defensiveness:
a metric element without a target is skipped rather than animated to \code{NaN}. Its
threshold is 0.5, so the numbers start counting only when the band is half visible.

Note the callback is passed as a bare function name, \code{animateCounter}, rather
than wrapped in another arrow function. That is cleaner than the sibling, which
writes \code{(el) => \{ animateCounter(el); \}} for the same effect.

\begin{honest}{The class \code{counted} is added and styled nowhere}
The third call passes \code{'counted'} rather than \code{'visible'}, so each metric
number gains a class no CSS rule and no other JavaScript ever reads. It is harmless
and could serve as a marker for a developer inspecting the page, but nothing uses
it. Note also that the metric numbers are children of \code{.metrics-grid}, which is
a \code{.stagger}, so they are already faded in by the second call: this class is
genuinely spare.
\end{honest}

\paragraph{Lines 53 to 64: smooth anchor scrolling, and the close.}

\begin{lstlisting}[language=Java,firstnumber=53,caption={\file{js/main.js}, lines 53--64}]

  /* ── Smooth anchor scroll ── */
  document.querySelectorAll('a[href^="#"]').forEach((a) => {
    a.addEventListener('click', (e) => {
      const t = document.querySelector(a.getAttribute('href'));
      if (!t) return;
      e.preventDefault();
      t.scrollIntoView({ behavior: 'smooth', block: 'start' });
    });
  });

})();
\end{lstlisting}

Line 55 finds every link whose destination starts with a hash.
\code{a[href\^{}="\#"]} is an attribute selector and \code{\^{}=} means ``starts
with''. On this page that matches three links: the navigation's ``Book a call'', the
hero's ``Book a call'', and the hero's ``See the work''.

Line 57 looks up the destination using the link's own \code{href} as a selector,
which works because \code{\#book} is valid CSS meaning ``the element with id book''.

Line 58 returns early if the destination does not exist, leaving the browser to
handle the click normally. Line 59 then cancels the browser's default jump, but only
after the destination is confirmed, which is the correct order: cancel first and a
broken link would do nothing at all instead of at least trying.

Line 60 performs the scroll, with \code{block: 'start'} aligning the destination to
the top of the window.

\begin{honest}{This whole block is redundant, and it carries a latent crash}
\file{css/site.css} line 28 already sets \code{scroll-behavior: smooth} on the
document, making every anchor link glide with no JavaScript at all. These ten lines
reproduce that.

They also introduce a failure the CSS version does not have. A link written
\code{href="\#"}, a common way to write a placeholder link, would reach line 57 as
\code{document.querySelector('\#')}. That is not a valid selector, so the browser
throws a \code{SyntaxError}, and because it happens inside a click handler it would
break that click and leave an error in the console.

No such link exists today, so this is latent rather than live. The honest summary is
that the safest change to this file would be to delete lines 53 to 62 entirely: it
would remove a redundancy and a hazard in one edit, and nothing visible would
change.
\end{honest}

Line 64 closes and immediately runs the wrapper opened on line 4.

\subsubsection*{What it connects to}

\textbf{Named by:} \file{new-design-reference/index.html}, line 393, and nothing
else.

\textbf{Names:} no file at all. It imports nothing and fetches nothing.

\textbf{Its entire interface with the rest of the project is six strings}, worth
listing because they are the contract nothing checks:

\begin{table}[H]\centering\small
\begin{tabular}{@{}lll@{}}
\toprule
String & Used at & Where it must also exist \\
\midrule
\code{.nav}        & line 8  & \file{index.html} line 23, \file{site.css} line 91 \\
\code{scrolled}    & line 10 & \file{site.css} line 101 \\
\code{.reveal}     & line 50 & \file{index.html} (five places), \file{site.css} line 75 \\
\code{.stagger}    & line 51 & \file{index.html} (two places), \file{site.css} line 82 \\
\code{visible}     & lines 50--51 & \file{site.css} lines 80 and 88 \\
\code{.metric-num} & line 52 & \file{index.html} (four places), \file{site.css} line 407 \\
\bottomrule
\end{tabular}
\caption{Every point of contact between this file and the other two. Break any one
string in one place and the page keeps working, minus one behaviour, with no error
reported anywhere.}
\end{table}

\subsubsection*{How it breaks}

\begin{itemize}
\item \textbf{Load it in the \code{<head>} instead of at the end of the body} and
  all three behaviours vanish silently, since every lookup runs immediately and
  finds nothing.
\item \textbf{Rename a class in the HTML} and the corresponding behaviour stops,
  silently, because of the early-return guards.
\item \textbf{Delete this file} and the page is static, with four zeros in the
  metrics band. The hero statistics still read correctly, which limits the damage.
\item \textbf{Add a link with \code{href="\#"}} and clicking it throws a
  \code{SyntaxError} from line 57.
\item \textbf{A visitor who has asked for reduced motion} still gets the counting
  animation, because this file never checks.
\item \textbf{An ancient browser without \code{IntersectionObserver}} would throw on
  line 39 and lose all three reveal behaviours, which would leave every
  \code{.reveal} element permanently invisible were it not for the CSS
  reduced-motion block. Every browser in current use supports it, so this is a
  historical note rather than a live concern.
\end{itemize}

\subsubsection*{Honest findings for this file}

The counter ignores \code{prefers-reduced-motion}; the smooth-scroll block
duplicates a CSS feature and carries a latent \code{SyntaxError}; the
\code{counted} class is inert; and the \code{||} defaults on lines 42 and 46 cannot
express a falsy argument. Against those: the file is genuinely well written, with
consistent guards, a sensible shared helper, correct use of \code{unobserve} and
\code{passive}, no dependencies, and slightly tighter than its sibling in two
places.
